\section{Rational equivalence}
\label{section-rational-equivalence}

\noindent
In this section we define {\it rational equivalence} on $k$-cycles.
We will allow locally finite sums of images of
principal divisors (under closed immersions). This leads to some
pretty strange phenomena, see Example \ref{example-weird}.
However, if we do not allow these then we do not know how to prove that
capping with Chern classes of line bundles factors through rational
equivalence.

\begin{definition}
\label{definition-rational-equivalence}
Let $(S, \delta)$ be as in Situation \ref{situation-setup}.
Let $X$ be a scheme locally of finite type over $S$.
Let $k \in \mathbf{Z}$.
\begin{enumerate}
\item Given any locally finite collection $\{W_j \subset X\}$
of integral closed subschemes with $\dim_\delta(W_j) = k + 1$,
and any $f_j \in R(W_j)^*$ we may consider
$$
\sum (i_j)_*\text{div}(f_j) \in Z_k(X)
$$
where $i_j : W_j \to X$ is the inclusion morphism.
This makes sense as the morphism
$\coprod i_j : \coprod W_j \to X$ is proper.
\item We say that $\alpha \in Z_k(X)$ is {\it rationally equivalent to zero}
if $\alpha$ is a cycle of the form displayed above.
\item We say $\alpha, \beta \in Z_k(X)$ are
{\it rationally equivalent} and we write $\alpha \sim_{rat} \beta$
if $\alpha - \beta$ is rationally equivalent to zero.
\item We define
$$
\CH_k(X) = Z_k(X) / \sim_{rat}
$$
to be the {\it Chow group of $k$-cycles on $X$}. This is sometimes called
the {\it Chow group of $k$-cycles modulo rational equivalence on $X$}.
\end{enumerate}
\end{definition}

\noindent
There are many other interesting (adequate) equivalence relations.
Rational equivalence is the coarsest one of them all.

\begin{remark}
\label{remark-chow-group-pointwise}
Let $(S, \delta)$ be as in Situation \ref{situation-setup}.
Let $X$ be locally of finite type over $S$. Let $k \in \mathbf{Z}$.
Let us show that we have a presentation
$$
\bigoplus\nolimits_{\delta(x) = k + 1}' K_1^M(\kappa(x))
\xrightarrow{\partial}
\bigoplus\nolimits_{\delta(x) = k}' K_0^M(\kappa(x)) \to
\CH_k(X) \to 0
$$
Here we use the notation and conventions introduced in
Remark \ref{remark-cycles-pointwise} and in addition
\begin{enumerate}
\item $K_1^M(\kappa(x)) = \kappa(x)^*$ is the degree $1$ part of
the Milnor K-theory of the residue field $\kappa(x)$ of the point
$x \in X$ (see Remark \ref{remark-gersten-complex-milnor}), and
\item the differential $\partial$ is defined as follows:
given an element $\xi = \sum_x f_x$ we denote $W_x = \overline{x}$
the integral closed subscheme of $X$ with generic point $x$ and we set
$$
\partial(\xi) = \sum (W_x \to X)_*\text{div}(f_x)
$$
in $Z_k(X)$ which makes sense as we have seen that the second
term of the complex is equal to $Z_k(X)$ by
Remark \ref{remark-cycles-pointwise}.
\end{enumerate}
The fact that we obtain a presentation of $\CH_k(X)$ follows
immediately by comparing with Definition \ref{definition-rational-equivalence}.
\end{remark}

\noindent
A very simple but important lemma is the following.

\begin{lemma}
\label{lemma-restrict-to-open}
Let $(S, \delta)$ be as in Situation \ref{situation-setup}.
Let $X$ be a scheme locally of finite type over $S$.
Let $U \subset X$ be an open subscheme, and denote
$i : Y = X \setminus U \to X$ as a reduced closed subscheme of $X$.
Let $k \in \mathbf{Z}$.
Suppose $\alpha, \beta \in Z_k(X)$.
If $\alpha|_U \sim_{rat} \beta|_U$ then there exist a cycle
$\gamma \in Z_k(Y)$ such that
$$
\alpha \sim_{rat} \beta + i_*\gamma.
$$
In other words, the sequence
$$
\xymatrix{
\CH_k(Y) \ar[r]^{i_*} & \CH_k(X) \ar[r]^{j^*} & \CH_k(U) \ar[r] & 0
}
$$
is an exact complex of abelian groups.
\end{lemma}

\begin{proof}
Let $\{W_j\}_{j \in J}$ be a locally finite collection of integral closed
subschemes of $U$ of $\delta$-dimension $k + 1$, and let $f_j \in R(W_j)^*$
be elements such that $(\alpha - \beta)|_U = \sum (i_j)_*\text{div}(f_j)$
as in the definition. Set $W_j' \subset X$ equal
to the closure of $W_j$. Suppose that $V \subset X$ is a quasi-compact
open. Then also $V \cap U$ is quasi-compact open in $U$ as
$V$ is Noetherian. Hence the set
$\{j \in J \mid W_j \cap V \not = \emptyset\}
= \{j \in J \mid W'_j \cap V \not = \emptyset\}$
is finite since $\{W_j\}$ is locally finite. In other words we see that
$\{W'_j\}$ is also locally finite. Since $R(W_j) = R(W'_j)$ we see
that
$$
\alpha - \beta - \sum (i'_j)_*\text{div}(f_j)
$$
is a cycle supported on $Y$ and the lemma follows (see
Lemma \ref{lemma-exact-sequence-open}).
\end{proof}

\begin{lemma}
\label{lemma-exact-sequence-closed-chow}
Let $(S, \delta)$ be as in Situation \ref{situation-setup}.
Let $X$ be locally of finite type over $S$. Let $X_1, X_2 \subset X$
be closed subschemes such that $X = X_1 \cup X_2$ set theoretically.
For every $k \in \mathbf{Z}$ the sequence of abelian groups
$$
\xymatrix{
\CH_k(X_1 \cap X_2) \ar[r] &
\CH_k(X_1) \oplus \CH_k(X_2) \ar[r] &
\CH_k(X) \ar[r] &
0
}
$$
is exact. Here $X_1 \cap X_2$ is the scheme theoretic intersection and the
maps are the pushforward maps with one multiplied by $-1$.
\end{lemma}

\begin{proof}
By Lemma \ref{lemma-exact-sequence-closed} the arrow
$\CH_k(X_1) \oplus \CH_k(X_2) \to \CH_k(X)$ is surjective.
Suppose that $(\alpha_1, \alpha_2)$ maps to zero under this map.
Write $\alpha_1 = \sum n_{1, i}[W_{1, i}]$ and
$\alpha_2 = \sum n_{2, i}[W_{2, i}]$. Then we obtain a locally
finite collection $\{W_j\}_{j \in J}$ of integral closed
subschemes of $X$ of $\delta$-dimension $k + 1$ and $f_j \in R(W_j)^*$
such that
$$
\sum n_{1, i}[W_{1, i}] + \sum n_{2, i}[W_{2, i}] = \sum (i_j)_*\text{div}(f_j)
$$
as cycles on $X$ where $i_j : W_j \to X$ is the inclusion morphism.
Choose a disjoint union decomposition $J = J_1 \amalg J_2$ such that
$W_j \subset X_1$ if $j \in J_1$ and $W_j \subset X_2$ if $j \in J_2$.
(This is possible because the $W_j$ are integral.) Then we can write
the equation above as
$$
\sum n_{1, i}[W_{1, i}] - \sum\nolimits_{j \in J_1} (i_j)_*\text{div}(f_j) =
- \sum n_{2, i}[W_{2, i}] + \sum\nolimits_{j \in J_2} (i_j)_*\text{div}(f_j)
$$
Hence this expression is a cycle (!) on $X_1 \cap X_2$. In other words
the element $(\alpha_1, \alpha_2)$ is in the image of the first arrow
and the proof is complete.
\end{proof}

\begin{example}
\label{example-weird}
Here is a ``strange'' example.
Suppose that $S$ is the spectrum of a field $k$
with $\delta$ as in Example \ref{example-field}.
Suppose that $X = C_1 \cup C_2 \cup \ldots$ is an infinite
union of curves $C_j \cong \mathbf{P}^1_k$ glued together
in the following way: The point $\infty \in C_j$ is glued
transversally to the point $0 \in C_{j + 1}$ for $j = 1, 2, 3, \ldots$.
Take the point $0 \in C_1$. This gives a zero cycle
$[0] \in Z_0(X)$. The ``strangeness'' in this situation is
that actually $[0] \sim_{rat} 0$! Namely we can choose
the rational function $f_j \in R(C_j)$ to be the function
which has a simple zero at $0$ and a simple pole at $\infty$
and no other zeros or poles. Then we see that the sum
$\sum (i_j)_*\text{div}(f_j)$ is exactly the $0$-cycle
$[0]$. In fact it turns out that $\CH_0(X) = 0$ in this example.
If you find this too bizarre, then you can just
make sure your spaces are always quasi-compact
(so $X$ does not even exist for you).
\end{example}

\begin{remark}
\label{remark-infinite-sums-rational-equivalences}
Let $(S, \delta)$ be as in Situation \ref{situation-setup}.
Let $X$ be a scheme locally of finite type over $S$.
Suppose we have infinite collections $\alpha_i, \beta_i \in Z_k(X)$,
$i \in I$ of $k$-cycles on $X$. Suppose that the supports
of $\alpha_i$ and $\beta_i$ form locally finite collections
of closed subsets of $X$ so that $\sum \alpha_i$
and $\sum \beta_i$ are defined as cycles. Moreover, assume that
$\alpha_i \sim_{rat} \beta_i$ for each $i$. Then it is not
clear that $\sum \alpha_i \sim_{rat} \sum \beta_i$. Namely,
the problem is that the rational equivalences may be
given by locally finite
families $\{W_{i, j}, f_{i, j} \in R(W_{i, j})^*\}_{j \in J_i}$
but the union $\{W_{i, j}\}_{i \in I, j\in J_i}$ may not
be locally finite.

\medskip\noindent
In many cases in practice, one has a locally finite family of closed
subsets $\{T_i\}_{i \in I}$ such that $\alpha_i, \beta_i$
are supported on $T_i$ and such that $\alpha_i = \beta_i$
in $\CH_k(T_i)$, in other words, the families
$\{W_{i, j}, f_{i, j} \in R(W_{i, j})^*\}_{j \in J_i}$
consist of subschemes $W_{i, j} \subset T_i$. In this case it is true that
$\sum \alpha_i \sim_{rat} \sum \beta_i$ on $X$, simply because
the family $\{W_{i, j}\}_{i \in I, j\in J_i}$ is automatically
locally finite in this case.
\end{remark}






